This is a small CPU study. (i) One algorithm, one graph generator and two test families; conclusions may not transfer to other CLRS algorithms, to integer weights, or to directed graphs. (ii) Five seeds for the main table and three for the ablations; the seed-to-seed spread is as large as most between-system differences, so most comparisons at 32 and 64 nodes are underpowered and only one registered test (MLP at S16) is significant. Medians, the failure threshold and the failure counts are descriptive choices made after seeing the data, and the post hoc tests are uncorrected. (iii) A single training budget of 1000 steps with no hyperparameter search for any system; the negative result for step supervision may reverse with longer training, different loss weights, teacher forcing or a different processor, and the MLP may be under-trained. (iv) The architectures are our simplified reimplementations (width 32, no gating, no layer norm, no triplet reasoning) of the cited ideas, not the published systems, so they say nothing about the published results. (v) Absolute errors are not comparable across the two families because distance scales differ; we ran no trivial reference predictor (such as a constant), so the only anchor for the relative error is that predicting zero gives 1. (vi) The training graphs are small enough that the training distribution is only partly representative (very few distinct graphs at $n=4$). (vii) The step sweep covers only the two step-supervised systems on three seeds; whether the final-only models tolerate extra steps is unknown. (viii) The mechanism behind the failures (accumulating error, activation growth with degree, non-idempotent updates) is not isolated by any run in this paper, and the runs with non-finite outputs were not inspected beyond their final metrics. A full study would add more seeds, tuned baselines, larger training sizes, other algorithms and the published processors.
